\cleardoublepage
\section{Summary and Tips}

\subsection*{Interesting Facts}

\begin{itemize}
  \item Ratios are used to mix paint colors, create scale models, adjust ingredients in recipes, and describe odds in games!
  \item Rates are everywhere—in recipes, gas mileage (MPG), data usage on your phone, and how fast athletes run.
  \item Unit rates are super helpful when budgeting, comparing phone plans, or deciding which product gives you the best deal!
  \item Proportions are used in architecture, photography, and fashion design to maintain balance and scale!
  \item The symbol \textbf{\%} was first used hundreds of years ago in Europe. Percentages are now used globally in business, science, and even sports statistics!
\end{itemize}

\subsection*{Tips for Students}

\subsubsection*{Ratios}
\begin{itemize}
  \item Keep the order of comparison the same—read carefully.
  \item Simplify ratios when possible, just like you simplify fractions.
  \item Think visually—use pictures or charts if it helps you understand.
\end{itemize}

\subsubsection*{Rates}
\begin{itemize}
  \item Always include the units (like miles, hours, or dollars).
  \item To find a rate, divide the first number by the second.
  \item Use rates to compare which deal is better when shopping.
\end{itemize}

\subsubsection*{Unit Rate}
\begin{itemize}
  \item Make sure the second number becomes 1 (that's what "unit" means!).
  \item Use real-world examples when practicing—like prices, distance, or time.
  \item Always include the unit in your answer (like "per hour" or "per item").
\end{itemize}

\subsubsection*{Proportions}
\begin{itemize}
  \item Always simplify ratios before comparing.
  \item Use cross multiplication to solve quickly.
  \item Write word problems as proportions to organize your thinking.
\end{itemize}

\subsubsection*{Percent}
\begin{itemize}
  \item Think of \textbf{\%} as "out of 100."
  \item Use visual aids like 100 grids or pie charts to understand percent.
  \item Practice converting between fractions, decimals, and percents until it feels natural.
\end{itemize}











%%%%%%%%%%%%%%%%%%%%%%%%%%%%%%%%%%%%%%%%%%%%%%%%%%%%%%%
%%%%%%%%%%%%%%%%%%%%%%%%%%%%%%%%%%%%%%%%%%%%%%%%%%%%%%%
%%%%%%%%%%%%%%%%%%%%%%%%%%%%%%%%%%%%%%%%%%%%%%%%%%%%%%%
%%%%%%%%%%%%%%%%%%%%%%%%%%%%%%%%%%%%%%%%%%%%%%%%%%%%%%%
%%%%%%%%%%%%%%%%%%%%%%%%%%%%%%%%%%%%%%%%%%%%%%%%%%%%%%%
%%%%%%%%%%%%%%%%%%%%%%%%%%%%%%%%%%%%%%%%%%%%%%%%%%%%%%%
\newpage
\begin{itemize}
  \item The word \inlinequote{digit \say{digit2} is the one!} come ...
  \item The \inlinequote{word} ``digit'' come ... 
\end{itemize}




\blockcite{kim_ssat-sat_2003}{To be, or not to be, that is the question: 
Whether's tis nobler in the mind to suffer. The slings 
and arrows of outrageous fortune...}

\blockcite[359]{kim_ssat-sat_2003}{This is your quoted text.}


\subsection*{Tip for Students}

Example sentense is here. \parencite[p.~359]{kim_ssat-sat_2003}.


\say{Example sentense is here.} The author said this \parencite[p.~200]{noauthor_help_2010}.